% ============================================================
% CHAPTER 2 — Background and Theoretical Foundations (REVISED v2)
% ============================================================

\chapter{Background and Theoretical Foundations}
\label{chap:background}

This chapter provides the conceptual and formal foundations upon which the
thesis is built.  Section~\ref{sec:bpm} introduces Business Process Management
and its lifecycle.  Section~\ref{sec:declarative} surveys declarative process
modelling and presents the formal definition of DCR Graphs used throughout this
work.  Section~\ref{sec:rl} establishes the reinforcement-learning framework,
covering Markov Decision Processes, Proximal Policy Optimisation, and safe
reinforcement learning via shielding.
Section~\ref{sec:moopt} introduces multi-objective optimisation, Pareto
dominance, and linear scalarisation as a strategy for approximating the Pareto
front.  Finally, Section~\ref{sec:resources} discusses resource allocation in
BPM and its role in performance optimisation.

% ============================================================
\section{Business Process Management}
\label{sec:bpm}
% ============================================================

A \emph{business process} is a collection of events, activities, and decision
points that involve a number of actors and objects, collectively leading to an
outcome that has value to at least one
customer~\cite{Dumas2018FundamentalsManagement}.  \ac{BPM} is the discipline of discovering, analysing, redesigning, and
continuously monitoring such processes in order to achieve organisational
goals~\cite{Dumas2018FundamentalsManagement}.  Rather than optimising isolated
tasks, BPM is concerned with the end-to-end \emph{chain of events} that
constitutes a process instance, with common performance objectives including
cost reduction, shorter cycle times, and lower error
rates~\cite{Huang2011ReinforcementManagement}.

The management of business processes follows a continuous lifecycle comprising
process identification, discovery, analysis, redesign, implementation, and
monitoring~\cite{Dumas2018FundamentalsManagement}.  The present thesis is
primarily located in the \emph{analysis} and \emph{redesign} phases: it
provides an automated mechanism for identifying which among all compliant
process executions are most efficient with respect to cost and duration,
thereby supporting evidence-based redesign decisions.

Process models in BPM are classically divided into two paradigms.
\emph{Imperative} models, such as BPMN or Petri nets, prescribe an explicit
control flow specifying exactly how a process should proceed.  They excel at
describing rigid, well-structured workflows but struggle to represent flexible
or knowledge-intensive processes where rules may apply but sequencing is
discretionary~\cite{Groefsema2024SupportingFamilies}.  \emph{Declarative}
models, by contrast, specify only the constraints that must be respected,
allowing any execution that does not violate them; they are discussed in depth
in the next section.

% ============================================================
\section{Declarative Process Modelling}
\label{sec:declarative}
% ============================================================

\acp{DPM} represent the control flow of a process implicitly: they define a
set of constraints governing the execution of activities and allow any
execution that does not violate those
constraints~\cite{Trinh2024OnOfConstraints}.  This \emph{open-world assumption}
makes declarative models particularly well suited to knowledge-intensive
processes, where the sequence of activities depends on the context and expertise
of the actors involved~\cite{Groefsema2024SupportingFamilies}.

Among declarative notations, \emph{\ac{DCR}},
introduced by Hildebrandt and
Mukkamala~\cite{Hildebrandt2011DeclarativeGraphs}, have gained significant
traction in academia and in industrial practice,
notably through adoption in Danish public-sector case
management~\cite{Slaats2013Business2013} and the DCR Solutions
toolchain~\cite{Marquard2015Business2015,Christfort2025DCR-JS:Mining}.

\subsection{Formal Definition}

\begin{definition}[DCR Graph~{\cite{Diaz2024Pareto-OptimalModels}}]
\label{def:dcr}
A \emph{Dynamic Condition Response Graph} is a tuple
$G = \langle E, M, \mathit{Act}, \rightarrow\!\bullet,\; \bullet\!\rightarrow,\;
\pm,\; l \rangle$, where:
\begin{enumerate}
  \item $E$ is a finite set of \emph{events};
  \item $M \in \mathcal{M}(G) = \mathcal{P}(E) \times \mathcal{P}(E) \times
    \mathcal{P}(E)$ is a \emph{marking}, and $\mathcal{M}(G)$ is the set of
    all markings;
  \item $\mathit{Act}$ is the set of \emph{actions};
  \item $\rightarrow\!\bullet \subseteq E \times E$ is the \emph{condition}
    relation;
  \item $\bullet\!\rightarrow \subseteq E \times E$ is the \emph{response}
    relation;
  \item $\pm : E \times E \rightharpoonup \{+, \%\}$ defines dynamic
    \emph{inclusion} ($\rightarrow^+$) and \emph{exclusion}
    ($\rightarrow^\%$) relations;
  \item $l : E \rightarrow \mathit{Act}$ is a labelling function.
\end{enumerate}
\end{definition}

The marking $M = (\mathit{Ex}, \mathit{Re}, \mathit{In})$ encodes three sets of
events: \emph{executed} events~($\mathit{Ex}$), \emph{pending} responses
($\mathit{Re}$) that must eventually occur or be excluded, and currently
\emph{included} events~($\mathit{In}$).  An event~$e$ is \emph{enabled} in
marking~$M$ if (i)~$e \in \mathit{In}$ and (ii)~all its included condition
predecessors have been executed, i.e.,
$\{f \in \mathit{In} \mid f \rightarrow\!\bullet\, e\} \subseteq
\mathit{Ex}$~\cite{Diaz2024Pareto-OptimalModels}.

\subsection{Execution Semantics}

The execution semantics of a DCR graph is defined as a \emph{\ac{LTS}}~\cite{Diaz2024Pareto-OptimalModels}:

\begin{definition}[Transitions~{\cite{Diaz2024Pareto-OptimalModels}}]
\label{def:dcr_transitions}
For a DCR graph~$G$, the corresponding \ac{LTS} $\mathcal{T}(G)$ has transitions
$M' \xrightarrow{(e,a)} M''$ where:
\begin{enumerate}
  \item $M' = (\mathit{Ex}', \mathit{Re}', \mathit{In}')$ is the marking
    before the transition;
  \item $M'' = (\mathit{Ex}' \cup \{e\},\; \mathit{Re}'',\; \mathit{In}'')$
    is the marking after;
  \item $e \in \mathit{In}'$ and $l(e) = a$ (event is included and labelled
    $a$);
  \item $\{e' \in \mathit{In}' \mid e' \rightarrow\!\bullet\, e\} \subseteq
    \mathit{Ex}'$ (all conditions satisfied);
  \item $\mathit{In}'' = (\mathit{In}' \cup \{e' \mid e
    \rightarrow^+ e'\}) \setminus \{e' \mid e \rightarrow^\% e'\}$;
  \item $\mathit{Re}'' = (\mathit{Re}' \setminus \{e\}) \cup \{e' \mid e
    \bullet\!\rightarrow e'\}$.
\end{enumerate}
\end{definition}

A trace $t = (e_0,a_0),(e_1,a_1),\ldots$ is \emph{accepting} if no event
remains simultaneously included and pending without having been
executed~\cite{Diaz2024Pareto-OptimalModels}.  This acceptance condition
defines the set of valid, complete process executions.

\subsection{The Flat Conformance Space}
\label{sec:flat_conformance}

The open-world semantics of DCR graphs ensures that every accepting trace is
formally compliant by construction.  However, this property reveals a critical
limitation for optimisation: all accepting traces are treated as \emph{equally
valid}, regardless of their cost, duration, or business utility.  Burattin
et al.~\cite{AndreaBurattin2021Exploring_Conformance_Space} characterise this
as the \emph{flat conformance space}, noting that purely declarative compliance
checking provides no mechanism to rank or differentiate valid executions.  The
present thesis addresses this gap directly by learning a policy that
\emph{un-flattens} the conformance space --- identifying the Pareto-optimal
subset of accepting traces with respect to competing business objectives.

% ============================================================
\section{Reinforcement Learning}
\label{sec:rl}
% ============================================================

\ac{RL} is a computational framework in which an agent
learns to make sequential decisions by interacting with an environment and
optimising a cumulative reward
signal~\cite{Sutton2018ReinforcementIntroduction}.  In contrast to supervised
learning, the agent receives no labelled examples; it discovers effective
behaviour through trial and error, guided only by delayed feedback from the
environment.

\subsection{Markov Decision Processes}
\label{sec:mdp}

The standard formal model for RL is the \emph{Markov Decision
Process}~(MDP)~\cite{Huang2011ReinforcementManagement,Alshiekh2017SafeShielding}:

\begin{definition}[\ac{MDP}]
\label{def:mdp}
An MDP is a tuple $\mathcal{M} = (S, A, P, R, \gamma)$, where:
\begin{itemize}
  \item $S$ is a finite set of \emph{states};
  \item $A$ is a finite set of \emph{actions};
  \item $P : S \times A \times S \to [0,1]$ is a \emph{transition function},
    where $P(s' \mid s,a)$ is the probability of reaching state~$s'$ from
    state~$s$ by taking action~$a$;
  \item $R : S \times A \times S \to \mathbb{R}$ is a \emph{reward function};
  \item $\gamma \in [0,1)$ is a \emph{discount factor} controlling the
    influence of future rewards.
\end{itemize}
\end{definition}

The \emph{Markov property} states that the transition distribution depends only
on the current state and action, not on the history.  A \emph{policy}
$\pi : S \to \mathcal{P}(A)$ maps states to distributions over actions.  The
\emph{value function}
$V^\pi(s) = \mathbb{E}_\pi\!\left[\sum_{t=0}^{\infty} \gamma^t r_t \mid
s_0=s\right]$
measures the expected discounted return from state~$s$ under policy~$\pi$.  The
optimal policy~$\pi^*$ maximises~$V^\pi$ for all states
simultaneously~\cite{Sutton2018ReinforcementIntroduction}.

The reward function is the sole specification of the agent's objective, and its design strongly influences both convergence speed and the behaviour ultimately learned~\cite{Sutton2018ReinforcementIntroduction}. Beyond the task reward, it is common practice to add auxiliary shaping terms that encode intermediate progress or discourage undesirable behaviour~\cite{Mataric1994RewardLearning,NgPolicyShaping}. The reward used in this thesis follows this practice; its exact form is given in Section~\ref{sec:reward}.

\subsection{Proximal Policy Optimisation}
\label{sec:ppo}

Policy gradient methods directly optimise a parameterised policy $\pi_\theta$
by ascending the gradient of the expected return.  A key challenge is that
large gradient steps can destabilise training by moving the policy into
poorly-explored regions.  \emph{Trust Region Policy Optimisation}~(TRPO)~\cite{Schulman2015TrustOptimization}
addresses this by constraining the KL divergence between successive policies,
but requires expensive second-order optimisation.

\emph{\ac{PPO}}, introduced by Schulman et
al.~\cite{Schulman2017ProximalAlgorithms}, achieves similar stability
guarantees using only first-order gradient updates.  PPO defines a
\emph{clipped surrogate objective}:

\begin{equation}
\label{eq:ppo}
L^{\mathrm{CLIP}}(\theta) = \hat{\mathbb{E}}_t\!\left[
  \min\!\left(r_t(\theta)\,\hat{A}_t,\;
  \mathrm{clip}\!\left(r_t(\theta),\, 1-\varepsilon,\, 1+\varepsilon\right)
  \hat{A}_t\right)\right],
\end{equation}

where $r_t(\theta) = \pi_\theta(a_t \mid s_t) / \pi_{\theta_{\mathrm{old}}}
(a_t \mid s_t)$ is the probability ratio between the updated and old policy,
$\hat{A}_t$ is the \emph{Generalised Advantage Estimate}(GAE)~\cite{Schulman2016High-DimensionalEstimation}, and
$\varepsilon$ is a clipping hyperparameter (typically $0.2$).  The clipping
removes the incentive for the ratio to move outside
$[1-\varepsilon, 1+\varepsilon]$, acting as a conservative trust region
without requiring constrained
optimisation~\cite{Schulman2017ProximalAlgorithms}.

PPO alternates between collecting rollouts with the current policy and
performing $K$~epochs of minibatch gradient ascent on $L^{\mathrm{CLIP}}$.
Its combination of sample efficiency, simplicity, and robustness has made it
a widely adopted on-policy algorithm in applied deep
RL~\cite{Schulman2017ProximalAlgorithms,Towers2025Gymnasium:Environments}.

PPO is selected as the primary learning algorithm of this thesis; the
concrete rationale, tied to the discrete role-conditioned action space
and the shielding mechanism, is given in
Section~\ref{sec:ppo_algorithm}.


Alternative algorithms considered include Soft Actor-Critic~(SAC) and
Constrained Policy Optimisation~(CPO).  SAC is an off-policy algorithm
designed primarily for continuous action
spaces \cite{Haarnoja2018SoftActor};
CPO handles constraints natively but requires second-order optimisation and
enforces compliance through the optimisation criterion rather than the
structural, environment-level mechanism adopted in this thesis.  PPO is
therefore the most direct fit for the discrete, shielded environment of this
framework.

\subsection{Safe Reinforcement Learning and Shielding}
\label{sec:shielding_background}

Standard RL algorithms maximise cumulative reward without regard for safety
constraints.  \emph{Safe reinforcement learning} addresses this limitation by
incorporating safety considerations into the learning
process~\cite{Garcia2015ALearning}.  Two broad families of
approaches can be distinguished: methods that modify the \emph{optimisation
criterion}, such as constrained MDPs~\cite{AltmanCONSTRAINEDPROCESSES} and Lagrangian
relaxation~\cite{Achiam2017ConstrainedOptimization}, and methods that
modify the agent's \emph{exploration behaviour} by incorporating external
knowledge~\cite{Garcia2015ALearning}.

\emph{Shielded reinforcement learning}, introduced by Alshiekh et
al.~\cite{Alshiekh2017SafeShielding}, belongs to the second family.  A \emph{shield} is a
reactive component, synthesised from a safety specification and an abstraction
of the environment dynamics, that restricts the agent's action space at every
time step so that only safe actions remain available.  The shield satisfies a
\emph{minimum interference} property: it restricts the agent only when an
unrestricted action would violate the safety
specification~\cite{Alshiekh2017SafeShielding}.

Two placement strategies are distinguished.  In \emph{preemptive} shielding,
the shield is placed before the agent and provides the set of safe actions from
which the agent must choose.  In \emph{post-posed} shielding, the shield
monitors the agent's selected action and substitutes it with a safe alternative
only when necessary~\cite{Alshiekh2017SafeShielding}.  Formally, preemptive shielding
transforms the original MDP~$\mathcal{M}$ into a restricted
MDP~$\mathcal{M}'$ in which, for each state~$s$, the available action set is
reduced to $A'_s \subseteq A_s$ by removing all actions that could lead to a
violation of the safety specification.

This thesis adopts a post-posed shielding architecture, derived directly from
the DCR graph semantics, in which the agent samples over the full action
space rather than a restricted one; the concrete mechanism and design
rationale are presented in Chapter~\ref{chap:method}.

% ============================================================
\section{Multi-objective Optimisation}
\label{sec:moopt}

Many real-world problems involve more than one objective that must be optimised
simultaneously, and these objectives frequently conflict.  In BPM, cost and
execution duration are prototypical examples: assigning more experienced (and
expensive) resources reduces duration but increases cost.  \emph{Multi-objective
optimisation}~(MOO) provides a principled framework for reasoning about such
trade-offs~\cite{ArtemisDoumeniParetoAgentSimulator}.

\subsection{Pareto Dominance and the Pareto Front}

\begin{definition}[Pareto Dominance]
\label{def:pareto}
Let $\mathbf{f}(x) = (f_1(x), \ldots, f_k(x))$ be a vector of $k$ objective
functions to be minimised.  A solution $x$ \emph{dominates} $x'$, written
$x \prec x'$, if
\[
  \forall\, i \in \{1,\ldots,k\}:\; f_i(x) \le f_i(x')
  \quad \wedge \quad
  \exists\, j :\; f_j(x) < f_j(x').
\]
A solution $x^*$ is \emph{Pareto-optimal} if no feasible solution dominates
it.  The collection of all Pareto-optimal solutions is the \emph{Pareto
front}~$\mathcal{P}^*$.
\end{definition}

The Pareto front represents a set of optimal trade-offs: improving any
objective from a Pareto-optimal solution necessarily degrades at least one
other
objective~\cite{ArtemisDoumeniParetoAgentSimulator,KirchdorferLukasFromPolicies}.
Rather than yielding a single best answer, it provides decision-makers with a
\emph{menu} of non-dominated options from which to select according to their
operational priorities.

\subsection{Linear Scalarisation}
\label{sec:lit_scalarisation}

Linear scalarisation is the most widely used technique for reducing a
multi-objective optimisation problem to a sequence of single-objective
sub-problems.  Given a weight vector $\mathbf{w} = (\alpha, \beta)$ with
$\alpha, \beta \ge 0$, the scalarised reward is defined as
\begin{equation}
\label{eq:scalar}
r_{\mathbf{w}} = \alpha \cdot f_1 + \beta \cdot f_2,
\end{equation}
where $f_1$ and $f_2$ are the objective functions to be optimised.  By
training independent policies under different weight vectors and aggregating
the resulting solutions, it is possible to approximate the Pareto front of a
multi-objective problem~\cite{Roijers2013ADecision-Making}.

The theoretical foundations and limitations of this approach are well
established.  Roijers et al.~\cite{Roijers2013ADecision-Making} formalise the
\emph{multi-objective MDP}~(MOMDP) and show that linear scalarisation can
recover only policies that lie on the \emph{convex hull} of the Pareto front
--- the \emph{Convex Coverage Set}~(CCS) --- and provably fails to recover
solutions in non-convex (concave) regions.  This limitation, first
demonstrated empirically in RL by Vamplew et
al.~\cite{Vamplew2008OnFronts}, has been confirmed across multiple domains
and architectures~\cite{Hayes2022APlanning}.

Despite this limitation, linear scalarisation remains the dominant baseline in
multi-objective RL for two reasons.  First, it is computationally tractable:
each weight configuration reduces to a standard single-objective RL problem,
so any existing algorithm (Q-learning, PPO, \ac{A3C}) can be applied without
modification.  Second, for processes with \emph{convex} Pareto fronts ---
which includes many practical resource allocation and scheduling problems ---
it recovers the full front exactly.  Hayes et al.~\cite{Hayes2022APlanning}
provide a practical guide to identifying when linear scalarisation is
sufficient and when non-linear alternatives such as Chebyshev scalarisation
are required.

The selection of weight vectors is itself a non-trivial design problem.  Das
and Dennis~\cite{DasDr.Das-NBIPaper-1998} propose the \emph{simplex-lattice
design}, which generates $\binom{H+m-1}{m-1}$ evenly-spaced weight vectors
for $m$ objectives and $H$ divisions per axis; for two objectives this yields
$H+1$ vectors from $(1,0)$ to $(0,1)$ in steps of $1/H$.  This construction
ensures even coverage of the weight simplex but does not guarantee even
coverage of objective space, since the mapping from weights to Pareto points
is generally non-isomorphic~\cite{Vamplew2008OnFronts}.

In this thesis, linear scalarisation with a simplex-lattice weight grid is
adopted as the primary strategy for Pareto front approximation.  The choice is
justified by the structure of the problem: cost and duration objectives in a
DCR process are bounded and exhibit a convex trade-off frontier, as confirmed
empirically in Chapter~\ref{chap:results}.

% ============================================================
\section{Resource Allocation in BPM}
\label{sec:resources}
% ============================================================

In a BPM context, a \emph{resource} is any actor or agent that carries out
process activities, typically a human worker, but potentially also a machine
or software system~\cite{Huang2011ReinforcementManagement}.  \emph{Resource
allocation} refers to the assignment of resources to work items at runtime, and
it is recognised as one of the principal levers for improving process
performance~\cite{Huang2011ReinforcementManagement}.

Resources may differ along two key dimensions: \emph{cost} and
\emph{capability}.  A more experienced (Expert) resource typically completes
an activity faster but at higher cost, while a junior resource incurs lower
cost at the price of longer execution time.  This cost--duration trade-off is
the primary source of multi-objective tension in this thesis and motivates the
design of role-conditioned action spaces in the proposed framework
(Chapter~\ref{chap:method}).

Resource allocation decisions can be modelled at different levels of
abstraction.  At the \emph{global} level, a centralised policy decides, for
each arriving work item, which available resource to assign, optimising an
aggregate objective such as total throughput or average cycle
time~\cite{Huang2011ReinforcementManagement}.  At the \emph{event} level ---
the granularity adopted in this thesis --- the decision is made for each
individual activity within a specific process instance, and both the process
state (DCR marking) and the resource profile are part of the decision context.
The relationship between these approaches and prior work is discussed in
Section~\ref{sec:lit_resources}.
